% !TEX program = xelatex
\documentclass[11pt,a4paper]{article}
\usepackage{algo-preamble}

\title{\textbf{L06 \textperiodcentered{} Γραφήματα I}\\[2pt]
\large Βασικές Έννοιες}
\author{Algorithms Class Hub}
\date{\small Αλγόριθμοι και Πολυπλοκότητα (K17) \textperiodcentered{} ΕΚΠΑ}

\begin{document}
\maketitle

\begin{center}\itshape\small
Εισαγωγή στα γραφήματα: ορισμοί, αναπαράσταση (πίνακας vs λίστα γειτνίασης),
διαδρομές, κύκλοι, δένδρα, και η διατύπωση των προβλημάτων s--t συνεκτικότητας
και απόστασης.
\end{center}

\section{Τι θα δούμε}

Με αυτή τη διάλεξη ξεκινάει το \textbf{μεγαλύτερο} κομμάτι του μαθήματος:
αλγόριθμοι σε \textbf{γραφήματα}. Σύμφωνα με τη στατιστική των τελευταίων
εξεταστικών, οι αλγόριθμοι σε γραφήματα αντιπροσωπεύουν περίπου το
\textbf{30--35\%} του βαθμού. Όλα τα επόμενα 4 μαθήματα χτίζουν αυτό το θεμέλιο.

Εδώ είναι το προετοιμαστικό κεφάλαιο. Δεν θα δούμε αλγόριθμο ακόμα --- θα μάθουμε
\textbf{τι είναι} ένα γράφημα, \textbf{πώς το αποθηκεύουμε} στον υπολογιστή, και
\textbf{τι ερωτήσεις} θα του κάνουμε στα επόμενα κεφάλαια.

\begin{intuition}
\textbf{Διαισθητικά:} ένα γράφημα είναι ο πιο γενικός τρόπος να εκφράσουμε
\textbf{σχέσεις} μεταξύ αντικειμένων. Αν υπάρχει «αυτό συνδέεται με εκείνο»,
υπάρχει γράφημα. Από κοινωνικά δίκτυα μέχρι τη ροή δεδομένων ενός προγράμματος
--- όλα είναι γραφήματα. Αυτό κάνει τους αλγορίθμους που μαθαίνουμε εδώ
\textbf{παντού} εφαρμόσιμους.
\end{intuition}

\section{Τι είναι ένα γράφημα}

\textbf{Μη κατευθυνόμενο (ή απλό) γράφημα:} $G = (V, E)$ όπου:
\begin{itemize}
  \item $V$ $=$ σύνολο \textbf{κορυφών} (vertices ή nodes ή κόμβοι).
  \item $E$ $=$ σύνολο \textbf{ακμών} (edges) μεταξύ ζευγαριών κορυφών. Κάθε
  ακμή είναι μη διατεταγμένο ζευγάρι $\{u, v\}$.
\end{itemize}

\textbf{Παράμετροι μεγέθους:} $n(G) = |V|$ (πλήθος κορυφών) και
$m(G) = |E|$ (πλήθος ακμών). Ο \textbf{βαθμός (degree)} μιας κορυφής $v$ είναι το
πλήθος των ακμών στις οποίες ανήκει.

\subsection{Παράδειγμα}

Θεώρησε ένα γράφημα 8 κορυφών:
\begin{itemize}
  \item $V = \{1, 2, 3, 4, 5, 6, 7, 8\}$, $n(G) = 8$.
  \item $E = \{\{1,2\}, \{1,3\}, \{2,3\}, \{2,4\}, \{4,5\}, \{2,5\}, \{3,5\},
  \{3,7\}, \{7,8\}, \{3,8\}, \{5,6\}\}$, $m(G) = 11$.
  \item $\deg(1) = 2$, $\deg(2) = 4$, $\deg(3) = 5$, $\deg(6) = 1$.
\end{itemize}

\begin{keypoint}
\textbf{Λήμμα (handshaking):} $\sum_{v \in V} \deg(v) = 2 |E|$. Κάθε ακμή
«μετριέται» στους δύο βαθμούς των άκρων της. Πολύ χρήσιμο σε αναλύσεις.
\end{keypoint}

\section{Πραγματικά παραδείγματα γραφημάτων}

\begin{center}
\begin{tabular}{@{}l l l@{}}
\toprule
Γράφημα & Κορυφές & Ακμές \\
\midrule
Δίκτυα μεταφορών            & διασταυρώσεις    & αυτοκινητόδρομοι \\
Δίκτυα επικοινωνιών         & υπολογιστές      & καλώδια οπτικών ινών \\
Παγκόσμιος ιστός            & ιστοσελίδες      & υπερσύνδεσμοι \\
Κοινωνικά δίκτυα            & άνθρωποι         & σχέσεις \\
Τροφική αλυσίδα             & είδη οργανισμών  & θηρευτής $\to$ θήραμα \\
Συστήματα λογισμικού        & συναρτήσεις      & κλήσεις \\
Χρονοπρογραμματισμός        & εργασίες         & περιορισμοί προτεραιότητας \\
Κυκλώματα                   & πύλες            & καλώδια \\
Δίκτυο συγκοινωνιών (μετρό) & σταθμοί          & διαδοχικές στάσεις \\
\bottomrule
\end{tabular}
\end{center}

\begin{intuition}
\textbf{Σημείωση:} το ίδιο πρόβλημα μπορεί να μοντελοποιηθεί ως διαφορετικά
γραφήματα ανάλογα με την ερώτηση. Π.χ.\ στο δίκτυο συγκοινωνιών: αν με ενδιαφέρει
«πόσοι σταθμοί χωρίζουν δύο σημεία;», οι κορυφές είναι σταθμοί. Αν με ενδιαφέρει
«πόσες αλλαγές γραμμών χρειάζονται;», οι κορυφές είναι γραμμές και ακμές οι
σταθμοί κοινής αλλαγής.
\end{intuition}

\section{Αναπαράσταση γραφημάτων}

Έχουμε δύο επιλογές για να αποθηκεύσουμε ένα γράφημα στον υπολογιστή. Η επιλογή
επηρεάζει τις πολυπλοκότητες όλων των αλγορίθμων.

\subsection{Πίνακας γειτνίασης (adjacency matrix)}

Ένας πίνακας $A$ διαστάσεων $n \times n$ όπου:
\[
A_{ij} = \begin{cases} 1 & \text{αν } \{i, j\} \in E \\ 0 & \text{διαφορετικά}
\end{cases}
\]
Για μη κατευθυνόμενο γράφημα ο πίνακας είναι \textbf{συμμετρικός}
($A_{ij} = A_{ji}$). Για το παραπάνω γράφημα:
\[
A =
\begin{pmatrix}
0 & 1 & 1 & 0 & 0 & 0 & 0 & 0 \\
1 & 0 & 1 & 1 & 1 & 0 & 0 & 0 \\
1 & 1 & 0 & 0 & 1 & 0 & 1 & 1 \\
0 & 1 & 0 & 0 & 1 & 0 & 0 & 0 \\
0 & 1 & 1 & 1 & 0 & 1 & 0 & 0 \\
0 & 0 & 0 & 0 & 1 & 0 & 0 & 0 \\
0 & 0 & 1 & 0 & 0 & 0 & 0 & 1 \\
0 & 0 & 1 & 0 & 0 & 0 & 1 & 0
\end{pmatrix}
\]

\textbf{Ιδιότητες:}
\begin{itemize}
  \item Χώρος: $\Theta(n^2)$ --- \textbf{πυκνή} αναπαράσταση.
  \item Έλεγχος αν υπάρχει ακμή $\{i, j\}$: $\Theta(1)$. $\checkmark$
  \item Επίσκεψη \textbf{όλων} των ακμών μιας κορυφής: $\Theta(n)$ (σαρώνεις τη
  γραμμή).
  \item Επίσκεψη όλων των ακμών του γραφήματος: $\Theta(n^2)$.
\end{itemize}

\subsection{Λίστα γειτνίασης (adjacency list)}

Πίνακας $n$ λιστών, μία ανά κορυφή, που περιέχει τους γείτονές της.

\begin{codebox}
\begin{verbatim}
1 -> 2, 3
2 -> 1, 3, 4, 5
3 -> 1, 2, 5, 7, 8
4 -> 2, 5
5 -> 2, 3, 4, 6
6 -> 5
7 -> 3, 8
8 -> 3, 7
\end{verbatim}
\end{codebox}

\textbf{Ιδιότητες:}
\begin{itemize}
  \item Χώρος: $\Theta(n + m)$ --- \textbf{αραιή} αναπαράσταση.
  \item Έλεγχος αν υπάρχει ακμή $\{i, j\}$: $O(\deg(i))$ (σαρώνω τη λίστα).
  \item Επίσκεψη όλων των γειτόνων μιας κορυφής $u$: $\Theta(\deg(u))$.
  $\checkmark$
  \item Επίσκεψη όλων των ακμών του γραφήματος: $\Theta(n + m)$.
\end{itemize}

\subsection{Ποιο επιλέγουμε;}

\begin{keypoint}
\textbf{Κανόνας:}
\begin{itemize}
  \item \textbf{Πυκνά γραφήματα} ($m = \Theta(n^2)$, π.χ.\ πλήρες γράφημα):
  πίνακας γειτνίασης --- ίδιος χώρος, γρήγορο edge lookup.
  \item \textbf{Αραιά γραφήματα} ($m = O(n)$ ή $O(n \log n)$, που είναι το πιο
  συνηθισμένο στην πράξη): λίστα γειτνίασης --- γραμμικός χώρος, γρήγορη
  iteration γειτόνων.
\end{itemize}
\textbf{Στο μάθημα} οι περισσότεροι αλγόριθμοι αναλύονται ως $O(n + m)$ ή
$O((n+m) \log n)$ --- υπονοούν λίστα γειτνίασης. Όταν βλέπεις πολυπλοκότητα με
$m$, υπόθεσε λίστα γειτνίασης.
\end{keypoint}

\section{Διαδρομές και συνεκτικότητα}

\subsection{Διαδρομή (path)}

Μια \textbf{διαδρομή} σε γράφημα $G = (V, E)$ είναι μια ακολουθία κορυφών
\[
P = v_1, v_2, \dots, v_{k-1}, v_k
\]
με την ιδιότητα ότι κάθε διαδοχικό ζεύγος $\{v_i, v_{i+1}\}$ είναι ακμή. Το
\textbf{μήκος} της διαδρομής είναι $k - 1$ ακμές.

\textbf{Απλή διαδρομή:} όλες οι κορυφές είναι διαφορετικές (κανείς κόμβος δεν
επαναλαμβάνεται).

\begin{intuition}
Διαφορά διαδρομής vs απλής διαδρομής: σε μη απλή μπορεί να γυρνάω στον ίδιο
κόμβο. Στις περισσότερες αναλύσεις του μαθήματος ενδιαφερόμαστε για απλές ---
π.χ.\ όταν μιλάμε για \textbf{συντομότερη} διαδρομή, αν περνούσαμε δύο φορές από
τον ίδιο κόμβο, θα μπορούσαμε να αφαιρέσουμε τον κύκλο και να βρούμε μικρότερη.
\end{intuition}

\subsection{Συνεκτικό γράφημα}

Ένα μη κατευθυνόμενο γράφημα είναι \textbf{συνεκτικό} (connected) αν για
\textbf{κάθε} ζευγάρι κορυφών $u, v$ υπάρχει διαδρομή μεταξύ τους.

Αν δεν είναι συνεκτικό, χωρίζεται σε \textbf{συνεκτικές συνιστώσες} (connected
components) --- μέγιστα συνεκτικά υποσύνολα.

\section{Κύκλοι}

Ένας \textbf{κύκλος} (cycle) είναι μια διαδρομή $v_1, v_2, \dots, v_k$ με:
\begin{enumerate}
  \item $\{v_1, v_k\} \in E$ --- η αρχή και το τέλος συνδέονται με ακμή.
  \item $k \ge 3$ --- αποτρέπει τετριμμένους «κύκλους» μήκους 2.
  \item Όλες οι κορυφές διακεκριμένες --- δεν περνάει δύο φορές από τον ίδιο
  κόμβο.
\end{enumerate}
Δηλαδή: \textbf{απλή διαδρομή που κλείνει}.

\textbf{Παράδειγμα:} στο γράφημά μας, η διαδρομή $1 \to 2 \to 4 \to 5 \to 3 \to
1$ είναι κύκλος μήκους 5.

\section{Δένδρα}

Ένα γράφημα είναι \textbf{δένδρο} (tree) αν είναι ταυτόχρονα \textbf{συνεκτικό}
και \textbf{δεν περιέχει κύκλο} (ακυκλικό).

\begin{keypoint}
\textbf{Θεώρημα.} Έστω $G$ γράφημα με $n$ κορυφές. Οποιεσδήποτε \textbf{2} από
τις παρακάτω προτάσεις συνεπάγονται την \textbf{3η}:
\begin{enumerate}
  \item Το $G$ είναι συνεκτικό.
  \item Το $G$ δεν περιέχει κύκλο.
  \item Το $G$ έχει $n - 1$ ακμές.
\end{enumerate}
\textbf{Αν 1+2} $\leftrightarrow$ είναι δένδρο. \textbf{Αν 1+3} $\leftrightarrow$
είναι δένδρο. \textbf{Αν 2+3} $\leftrightarrow$ είναι δένδρο (γιατί ένα ακυκλικό
γράφημα με $n-1$ ακμές είναι σίγουρα συνεκτικό).
\end{keypoint}

\textbf{Συνέπεια --- βασική:} σε δένδρο με $n$ κορυφές υπάρχει \textbf{μοναδικό
μονοπάτι} μεταξύ κάθε ζεύγους κορυφών. Αν είχε δύο μονοπάτια θα είχε κύκλο.

\subsection{Δένδρα με ρίζα}

\textbf{Δένδρο με ρίζα (rooted tree).} Διαλέγουμε μία κορυφή $r$ ως
\textbf{ρίζα} και προσανατολίζουμε κάθε ακμή ώστε να απομακρύνεται από τη ρίζα.
Δημιουργεί \textbf{ιεραρχία}:
\begin{itemize}
  \item \textbf{Γονέας (parent)} της $v$: ο γείτονας της $v$ που είναι
  \textbf{πιο κοντά} στη ρίζα.
  \item \textbf{Παιδιά (children)} της $v$: οι γείτονες που είναι \textbf{πιο
  μακριά} από τη ρίζα.
  \item \textbf{Φύλλο (leaf):} κορυφή χωρίς παιδιά.
\end{itemize}

\textbf{Εφαρμογές:} φυλογενετικά δένδρα (εξελικτική ιστορία ειδών), δομές
αρχείων, DOM στο HTML, AST (Abstract Syntax Tree) ενός προγράμματος, και δένδρα
αναζήτησης / αναδρομής (όπως είδαμε νωρίτερα με τη mergesort).

\section{Τα προβλήματα που θα λύσουμε στα επόμενα μαθήματα}

Πριν δούμε αλγορίθμους, ας \textbf{διατυπώσουμε} τα προβλήματα που θα μας
απασχολήσουν.

\textbf{Πρόβλημα 1: s--t συνεκτικότητα.}
\begin{quoteblock}
Δεδομένων κορυφών $s$ και $t$, υπάρχει διαδρομή από την $s$ στην $t$;
\end{quoteblock}

\textbf{Πρόβλημα 2: s--t απόσταση.}
\begin{quoteblock}
Δεδομένων κορυφών $s$ και $t$, ποιο είναι το \textbf{μήκος} της συντομότερης
διαδρομής από $s$ σε $t$;
\end{quoteblock}

\textbf{Εφαρμογές:}
\begin{itemize}
  \item \textbf{Κοινωνικά δίκτυα:} ποια είναι η απόσταση μεταξύ δύο ατόμων;
  \item \textbf{Επίλυση λαβυρίνθου:} υπάρχει διαδρομή από αρχή προς έξοδο;
  \item \textbf{Six Degrees of Kevin Bacon:} πόσες ταινίες απέχει κάποιος
  ηθοποιός από τον Kevin Bacon;
  \item \textbf{Τηλεπικοινωνίες:} ελάχιστος αριθμός ενδιάμεσων κορυφών σε
  δίκτυο.
  \item \textbf{Routing protocols:} το OSPF στο διαδίκτυο χρησιμοποιεί
  αλγορίθμους shortest path.
\end{itemize}

\begin{keypoint}
\textbf{Αυτά} είναι τα προβλήματα που θα λυθούν στο επόμενο μάθημα με
\textbf{BFS} και \textbf{DFS}. Παρακάτω θα γενικεύσουμε σε \textbf{κατευθυνόμενα}
γραφήματα και θα δούμε \textbf{διμερότητα} και \textbf{ισχυρή συνεκτικότητα}, και
ακόμη παρακάτω θα προχωρήσουμε στους κλασικούς αλγορίθμους shortest path με βάρη
(Dijkstra, Bellman-Ford) και ελάχιστων δένδρων (Prim, Kruskal).
\end{keypoint}

\begin{recapbox}
\textbf{Τι κρατάμε από αυτή τη διάλεξη}
\begin{enumerate}
  \item \textbf{Γράφημα $= (V, E)$}. $n$ κορυφές, $m$ ακμές. Βαθμός κορυφής.
  Σύνολο όλων των βαθμών $= 2m$.
  \item \textbf{Δύο αναπαραστάσεις:} πίνακας γειτνίασης ($n^2$ χώρος, $O(1)$
  edge lookup) ή λίστα γειτνίασης ($n+m$ χώρος, $O(\deg)$ neighbours). Λίστα
  γειτνίασης για αραιά, πίνακας για πυκνά.
  \item \textbf{Έννοιες:} διαδρομή, απλή διαδρομή, κύκλος, συνεκτικότητα,
  συνεκτική συνιστώσα.
  \item \textbf{Δένδρα:} συνεκτικό $+$ ακυκλικό $+$ $n-1$ ακμές. Οποιεσδήποτε
  δύο από αυτές $\to$ η τρίτη.
  \item \textbf{Δένδρο με ρίζα:} ιεραρχία γονέα-παιδιού. Φυλογενετικά δένδρα,
  AST, DOM.
  \item \textbf{Προβλήματα προς λύση:} s--t συνεκτικότητα, s--t απόσταση. Στο
  επόμενο μάθημα θα δούμε \textbf{BFS} και \textbf{DFS} που τα λύνουν.
\end{enumerate}
\end{recapbox}

\lecturefooter
\end{document}
